% Pista ana-25s-q3c · Anàlisi
% Procedència: PAU Catalunya 2025 · Convocatòria de setembre · Sèrie 3 · Exercici 3
\documentclass[11pt,a4paper]{article}
\usepackage{pista}
\pistainfo{Catalunya 2025}{Convocatòria de setembre}{Sèrie 3}

\begin{document}

\section*{\textcolor{blauPista}{Exercici 3}}

\noindent Una empresa de calçat esportiu ha creat una sabatilla amb amortiment per a minimitzar les lesions per sesamoïditis. Els beneficis generats per la venda d'aquest producte, en milers d'euros, segueixen una funció de la forma $f(x) = ax^{3} + bx^{2} + cx$, on $x$ són els anys transcorreguts des que la sabatilla va sortir a la venda i $a$, $b$ i $c$ són constants reals.

\subsection*{\textit{c)} \normalfont Calculeu els valors de $a$, $b$ i $c$ sabent que el primer any es van obtenir el màxim de beneficis, amb un valor de 8\,000 euros, i que el segon any va haver-hi un punt d'inflexió en els beneficis. \hspace*{\fill}\textnormal{[1~punt]}}

\pista{En primer lloc, calcula $f'(x)$ i $f''(x)$: les necessitaràs per traduir les dades.

\medskip
A continuació, tradueix cada dada de l'enunciat en una equació. Per exemple, que el segon any hi hagi un punt d'inflexió vol dir $f''(2) = 0$. Fixa't que el màxim del primer any amaga dues condicions: una sobre $f'(1)$ i una altra sobre $f(1)$. Per a la de $f(1)$, tingues en compte que $f$ dona els beneficis en milers d'euros.

\medskip
Finalment, resol el sistema de tres equacions amb les incògnites $a$, $b$ i $c$.}

\end{document}
